\documentclass[../main]{subfiles}
\begin{document}
\chapter{Diagramas}
\begin{center}
\includegraphics[scale=1]{images/10/flujo.png}
\end{center}
\section{Diagrama de flujo}
\img{images/10/flujo2.png}{Diagrama de flujo}

\info{Diagrama de Flujo}{El diagrama de flujo o flujograma o diagrama de actividades es la representación gráfica de un algoritmo o proceso. Se utiliza en disciplinas como programación, economía, procesos industriales y psicología cognitiva.
}
\subsection{Tipos de Diagrama}

\subsubsection{UML}
En Lenguaje Unificado de Modelado (UML), es un diagrama de actividades que representa los flujos de trabajo paso a paso. Un diagrama de actividades muestra el flujo de control general.

\subsubsection{SysML}

En SysML el diagrama ha sido extendido para indicar flujos entre pasos que mueven elementos físicos por ejemplo combustible o energía por ejemplo presión. Los cambios adicionales permiten al diagrama soportar mejor flujos de comportamiento y datos continuos.

Estos diagramas utilizan símbolos con significados definidos que representan los pasos del algoritmo, y representan el flujo de ejecución mediante flechas que conectan los puntos de inicio y de fin del proceso.

\subsection{Normas de trabajo}
Las siguientes son acciones previas a la realización del diagrama de flujo:

\begin{enumerate}
\item Definir qué se espera obtener del diagrama de flujo.
\item Identificar quién lo empleará y cómo.
\item Establecer el nivel de detalle requerido.
\item Determinar los límites del proceso a describir.
\end{enumerate}
Los pasos a seguir para construir el diagrama de flujo son:
\begin{enumerate}
    
\item Establecer el alcance del proceso a describir. De esta manera quedará fijado el comienzo y el final del diagrama. Frecuentemente el comienzo es la salida del proceso previo y el final la entrada al proceso siguiente.
\item Identificar y listar las principales actividades/subprocesos que están incluidos en el proceso a describir y su orden cronológico.
\item Si el nivel de detalle definido incluye actividades menores, listarlas también.
\item Identificar y listar los puntos de decisión.
\item Construir el diagrama respetando la secuencia cronológica y asignando los correspondientes símbolos.
\item Prueba para ver errores

\end{enumerate}
\begin{center}
\includegraphics[scale=1.4]{images/10/simbologia-ansi-diagrama-flujo.jpg}
\end{center}


\actividad{ Investigar y responder las siguientes preguntas}

\begin{enumerate}
    \item ¿Qué es un diagrama de flujo?
    
    \item ¿Cuál es el propósito de utilizar un diagrama de flujo?
    
    \item Menciona al menos tres símbolos comunes utilizados en los diagramas de flujo.
    
    \item ¿Qué forma se utiliza para representar un proceso o acción en un diagrama de flujo?
    
    \item ¿Cuál es la forma estándar para representar una decisión en un diagrama de flujo?
    
    \item ¿Cómo se conectan los elementos en un diagrama de flujo para indicar la secuencia?
    
    \item ¿Qué símbolo se usa para representar la entrada o salida de datos en un diagrama de flujo?
    
    \item ¿Cuál es la dirección de lectura habitual en un diagrama de flujo?
    
    \item ¿Qué significa un rombo en un diagrama de flujo?
    
    \item ¿Cómo se representa el inicio o fin de un proceso en un diagrama de flujo?
\end{enumerate}


\actividad{ Para el Proyecto}
\begin{enumerate}
    \item Elaborar un diagrama de flujo para el proceso de producción del proyecto así como también que incluya controles de calidad.
    \item Elaborar un diagrama de flujo para el proceso de compras de materia prima 
    \item Elaborar un diagrama de flujo para el proceso de ventas.
\end{enumerate}
 

\end{document}
